\documentclass[11pt,reqno]{amsart}
\usepackage[T1]{fontenc}
\usepackage[utf8]{inputenc}
\usepackage{lmodern}
\usepackage{amsmath,amssymb,amsfonts,amsthm,mathtools,mathrsfs}
\usepackage{enumitem}
\usepackage{microtype}
\usepackage[colorlinks=true,linkcolor=blue,citecolor=blue,urlcolor=blue]{hyperref}
\usepackage[nameinlink,capitalize]{cleveref}
\allowdisplaybreaks
\setlength{\emergencystretch}{2em}

\newtheorem{theorem}{Theorem}[section]
\newtheorem{proposition}[theorem]{Proposition}
\newtheorem{lemma}[theorem]{Lemma}
\newtheorem{corollary}[theorem]{Corollary}
\newtheorem{claim}[theorem]{Claim}
\theoremstyle{definition}
\newtheorem{definition}[theorem]{Definition}
\newtheorem{assumption}[theorem]{Assumption}
\newtheorem{example}[theorem]{Example}
\theoremstyle{remark}
\newtheorem{remark}[theorem]{Remark}

\newcommand{\A}{\mathcal A}
\newcommand{\B}{\mathcal B}
\newcommand{\C}{\mathbb C}
\newcommand{\D}{\mathrm D}
\newcommand{\E}{\mathbb E}
\newcommand{\N}{\mathbb N}
\newcommand{\R}{\mathbb R}
\newcommand{\T}{\mathbb T}
\newcommand{\X}{\mathcal X}
\newcommand{\Y}{\mathcal Y}
\newcommand{\cL}{\mathscr L}
\newcommand{\cP}{\mathscr P}
\newcommand{\cR}{\mathscr R}
\newcommand{\Piop}{\Pi}
\newcommand{\Id}{\mathrm{Id}}
\newcommand{\Var}{\operatorname{Var}}
\newcommand{\Tot}{\operatorname{Tot}}
\newcommand{\supp}{\operatorname{supp}}
\newcommand{\norm}[1]{\left\lVert #1\right\rVert}
\newcommand{\abs}[1]{\left\lvert #1\right\rvert}
\newcommand{\pair}[2]{\left\langle #1,#2\right\rangle}

\title[Bilateral graph-current mixing]{Bilateral Graph-Current Mixing and Third-Order Response for Dynamical Systems with Moving Singularities}
\author{Qian Qi}
\address{Peking University, Beijing 100871, China}
\email{qiqian@pku.edu.cn}
\subjclass[2020]{Primary 37C30, 37D50; Secondary 37A25, 47A55, 60F17}
\keywords{linear response, moving singularities, graph currents, transfer operators, higher-order response, dispersing billiards}
\date{August 28, 2026}

\begin{document}

\begin{abstract}
We develop a response theory for parameterized dynamical systems whose
singularity sets move.  The parameter derivative of the transfer operator is
resolved into same-occurrence graph-current atoms.  Two one-sided estimates
are retained separately: a source-side estimate contracts the past time while
allowing forward graph growth, and a target-side estimate contracts the future
time while allowing backward graph growth.  A strict crossed-rate inequality
interpolates these estimates into a product tail in the two time variables.
The resulting absolute double summability is stable under finite difference
quotients and is strong enough to totalize the complete third source over
homogeneity, intersection, and atlas refinements.  This yields a continuous
third spectral jet on a graded Banach scale.

The abstract theorem is realized by an open four-branch pinball-cylinder
family with three moving physical seams and a nonzero saltus source.  We prove
uniform product summability through fourth parameter order, convergence of
finite difference arrays in $\ell^1(\N^2)$, and a complete third-order Kato
calculus.  We also construct a genuinely nonconjugate radial family of
finite-horizon specular Sinai billiards for which the completely assembled
invariant-projector source has response of every finite order; exact
cohomological twists have explicit all-order eigendata.  Finally, we prove
maximality statements showing why local geometry or any finite prefix of a
source atlas cannot imply the required global product tail.
\end{abstract}

\maketitle
\tableofcontents

\section{Introduction}

Parameter response is classical for uniformly hyperbolic smooth dynamics and
for transfer operators that vary differentiably on a fixed Banach space.  The
situation changes when a discontinuity, grazing set, or collision partition
moves with the parameter.  Differentiating a branch formula then produces
interior terms together with currents supported on moving faces and their
forward or backward descendants.  Such currents may regularize only after a
specific physical assembly, and their correlation response contains two
independent time indices.  These features are not captured by a one-time
spectral-gap estimate.

The obstruction is already visible at the level of a single atom.  Let $L_a$
be a transfer operator, $\Pi_a$ its rank-one invariant projection, and
$Q_a=L_a-\Pi_a$.  If $J_{a,\omega}$ is a current associated with one physical
occurrence, the relevant array is
\begin{equation}\label{eq:intro-array}
 A_{m,n}^{\omega}(a;g,f)
 =\pair{Q_a^nJ_{a,\omega}Q_a^m g}{f},
 \qquad m,n\ge0.
\end{equation}
An estimate of the form $\rho^m+\rho^n$ is useless for absolute response,
because its double sum diverges.  The first contribution of this paper is a
bilateral mechanism that turns two crossed one-sided estimates into
$\rho_*^{m+n}$.

Our second contribution concerns higher order.  A third derivative is not a
single current.  It is a totalized sum of regular terms, one-face currents,
intersection currents, endpoint terms, and refinement defects.  We introduce a
directed third-source construction and prove a joint Cauchy theorem in a CM2
norm.  A graded scale records the exact regularity loss of each parameter
letter, so every third-order resolvent word is genuinely composable.

The third contribution is actualization.  The four-branch model in
\cref{sec:fourbranch} is elementary enough for every estimate to be computed,
yet it has moving physical seams and a nonzero saltus source.  It therefore
separates the new mechanism from exact conjugacy.  In
\cref{sec:specular} we give a different actual channel: a nonconjugate radial
family of finite-horizon specular Sinai billiards.  There the invariant density
is explicit in normalized collision coordinates, and differentiating
invariance produces a complete physical assembly whose singular currents
cancel exactly.

The distinction between these two examples is important.  The four-branch
family verifies the full bilateral packet for arbitrary declared observables.
The specular radial family verifies an all-order invariant-projector channel
and exact cohomological twists.  It does not imply arbitrary noncoboundary
third response from finite-horizon geometry alone.  The maximality results of
\cref{sec:maximality} explain why such an implication is false without an
additional recovery or product-tail theorem.

Our framework is complementary to the perturbative spectral theory of
Keller--Liverani, the higher response theory for piecewise expanding maps, and
the anisotropic-space theory for dispersing billiards; see
\cite{BaladiSmania2008,KellerLiverani1999,DemersLiverani2008,
BaladiDemersLiverani2018}.  Moving scatterers and nonstationary billiards were
studied by Stenlund--Young--Zhang \cite{StenlundYoungZhang2013}; their memory
loss estimates motivate, but do not themselves supply, the two-time current
summability required here.

\subsection*{Main results}
The abstract product-tail theorem is \cref{thm:bilateral}.  Finite difference
quotients are treated in \cref{thm:fdq}.  The CM2-to-third-source theorem is
\cref{thm:u3}.  The open moving-seam realization is
\cref{thm:fourbranch}.  The specular radial result is
\cref{thm:sinai-u3}.  The maximality theorem is
\cref{thm:maximality}.

\subsection*{Organization}
\Cref{sec:framework} defines same-occurrence currents and the graded operator
setting.  \Cref{sec:bilateral} proves the product tail.
\Cref{sec:fdq} upgrades finite difference quotients to the complete two-time
topology.  \Cref{sec:u3} totalizes the third source.  The two actual families
appear in \cref{sec:fourbranch,sec:specular}.  The final section records the
sharp no-go boundary.

\section{Same-occurrence currents and graded response}\label{sec:framework}

Let $\A$ be a compact parameter space, possibly decomposed into finitely many
chambers.  For $a\in\A$, let $T_a:M\to M$ be a nonsingular piecewise $C^r$ map.
All collision or branch spaces are transported to one measurable reference
space.  We write $L_a$ for the transported Perron--Frobenius operator and
assume that $1$ is a simple isolated eigenvalue with invariant projection
$\Pi_a$.

\begin{definition}[Occurrence registry]\label{def:registry}
An occurrence label $\omega$ consists of a physical face, an incident side, an
owner branch, an oriented branch word, a homogeneity itinerary, a parameter
chamber, and an atlas chart.  Two labels are identified only when these data
agree after passage to a common refinement.  Orientation reversal changes the
sign but not the underlying UID.
\end{definition}

The parameter derivative is represented on a source module by
\begin{equation}\label{eq:source-decomp}
 \partial_aL_a=K_a^{\mathrm{reg}}+
 \sum_{\omega\in\Omega_a}J_{a,\omega}.
\end{equation}
The sum is initially formal.  Its meaning is established by the estimates
below.  Distinct UIDs are never canceled before a common-refinement identity
has been proved.

We use a finite graded scale
\begin{equation}\label{eq:ladder}
 X_R\hookrightarrow X_{R-1}\hookrightarrow\cdots\hookrightarrow X_0,
\end{equation}
with continuous dense embeddings.  The centered operator $Q_a=L_a-\Pi_a$
preserves each level and the $j$th parameter letter loses exactly $j$ levels:
\begin{equation}\label{eq:graded-letter}
 G_a^{(j)}:X_r\longrightarrow X_{r-j},
 \qquad r\ge j.
\end{equation}
The same notation includes the represented total source at order $j$.

\begin{assumption}[Uniform spectral and testing structure]\label{ass:spectral}
For each level used below, $Q_a$ is power bounded and exponentially contracting
on the centered subspace.  There is a test space $Y_r\subset X_r^*$ such that
all graph-current pairings are defined.  Reduced resolvents
\[
 R_{a,r}=(I-Q_a)^{-1}:X_r^0\to X_r^0
\]
are uniformly bounded and continuous in $a$.
\end{assumption}

The two-time atom associated with a centered $g$ and a test $f$ is
\begin{equation}\label{eq:atom}
 A_{m,n}^{\omega}(a;g,f)
 =\pair{Q_a^nJ_{a,\omega}Q_a^m g}{f}.
\end{equation}
We emphasize that $m$ measures the source-side history and $n$ the
post-insertion evolution.  They are controlled by different graph mechanisms.

\section{Bilateral crossed envelopes}\label{sec:bilateral}

\begin{assumption}[Bilateral current packet]\label{ass:bilateral}
There are constants $0<\alpha,\beta<1$, $\Gamma_+,\Gamma_-\ge1$, and
nonnegative occurrence constants $C_\omega^-,C_\omega^+$ such that, uniformly
in $a$, $m,n$, and unit vectors in the declared spaces,
\begin{align}
 \abs{A_{m,n}^{\omega}}
 &\le C_\omega^-\alpha^m\Gamma_+^n,
 \label{eq:reverse}\
 \abs{A_{m,n}^{\omega}}
 &\le C_\omega^+\Gamma_-^m\beta^n.
 \label{eq:forward}
\end{align}
The first estimate is the reverse or source-side recovery estimate; the second
is the forward or target-side estimate.
\end{assumption}

Set
\[
 A_0=\log(1/\alpha),\quad B_0=\log(1/\beta),\quad
 C_0=\log\Gamma_+,\quad D_0=\log\Gamma_-.
\]

\begin{theorem}[Bilateral product-tail criterion]\label{thm:bilateral}
Assume \cref{ass:bilateral} and
\begin{equation}\label{eq:rate-gap}
 A_0B_0>C_0D_0.
\end{equation}
Define
\begin{equation}\label{eq:lambda-star}
 \lambda_*=
 \frac{B_0+D_0}{A_0+B_0+C_0+D_0}
\end{equation}
and
\begin{equation}\label{eq:rho-star}
 \log\rho_*=\frac{C_0D_0-A_0B_0}
 {A_0+B_0+C_0+D_0}.
\end{equation}
Then $0<\rho_*<1$ and
\begin{equation}\label{eq:product-tail}
 \abs{A_{m,n}^{\omega}}
 \le (C_\omega^-)^{\lambda_*}(C_\omega^+)^{1-\lambda_*}
 \rho_*^{m+n}.
\end{equation}
If
\begin{equation}\label{eq:occ-sum}
 \sum_{\omega}
 (C_\omega^-)^{\lambda_*}(C_\omega^+)^{1-\lambda_*}<\infty,
\end{equation}
then
\begin{equation}\label{eq:cm2-sum}
 \sum_{\omega}\sum_{m,n\ge0}\abs{A_{m,n}^{\omega}}<\infty.
\end{equation}
The convergence is uniform on compact parameter chambers and on bounded sets
of inputs and tests.
\end{theorem}

\begin{proof}
For $0\le\lambda\le1$, geometric interpolation of
\eqref{eq:reverse}--\eqref{eq:forward} gives
\[
 \abs{A_{m,n}^{\omega}}
 \le (C_\omega^-)^\lambda(C_\omega^+)^{1-\lambda}
 (\alpha^\lambda\Gamma_-^{1-\lambda})^m
 (\Gamma_+^\lambda\beta^{1-\lambda})^n.
\]
The logarithms of the two time factors are
\[
 -\lambda A_0+(1-\lambda)D_0,
 \qquad
 \lambda C_0-(1-\lambda)B_0.
\]
They coincide at \eqref{eq:lambda-star}.  Their common value is
\eqref{eq:rho-star}, which is negative precisely under
\eqref{eq:rate-gap}.  This proves \eqref{eq:product-tail}.  Summing first in
$m,n$ yields the factor $(1-\rho_*)^{-2}$; then
\eqref{eq:occ-sum} gives \eqref{eq:cm2-sum}.
\end{proof}

\begin{remark}[Sharpness of the rate window]
The existence of some interpolation parameter for which both time factors are
strictly smaller than one is equivalent to $A_0B_0>C_0D_0$.  Equality is a
critical line: the best common factor equals one and absolute product
summability is not forced.
\end{remark}

\begin{corollary}[Regularity-loss route]\label{cor:loss-route}
Suppose $K_a:X_{r+\ell}\to X_r$, $\Pi_aK_a=0$, and
\[
 \norm{Q_a^m}_{X_{r+\ell}^0\to X_{r+\ell}^0}
 \le C_+\rho_+^m,
 \qquad
 \norm{Q_a^n}_{X_r^0\to X_r^0}
 \le C_-\rho_-^n.
\]
Then
\[
 \abs{\pair{Q_a^nK_aQ_a^m g}{f}}
 \le C_-C_+\norm{K_a}\rho_-^n\rho_+^m
 \norm{g}_{X_{r+\ell}}\norm{f}_{X_r^*}.
\]
Thus a finite loss of regularity at the singular insertion is compatible with
CM2.
\end{corollary}

\section{Finite difference quotients in the complete topology}\label{sec:fdq}

For $h\ne0$ put
\[
 K_{a,h}=\frac{L_{a+h}-L_a}{h}.
\]
The centered evolutions on the two sides need not be identical, so the finite
array is
\[
 A_{m,n}^{h}(a;g,f)=
 \pair{Q_{a+h}^nK_{a,h}Q_a^m g}{f}.
\]

\begin{definition}
We say that finite difference quotients converge in the CM2 topology if
\begin{equation}\label{eq:fdq-def}
 \sum_{m,n\ge0}
 \abs{A_{m,n}^{h}(a;g,f)-A_{m,n}^{0}(a;g,f)}\longrightarrow0
\end{equation}
uniformly for $a$ in compact chambers and for $g,f$ in the declared unit
balls.
\end{definition}

\begin{theorem}[Uniform head--tail criterion]\label{thm:fdq}
Assume that:
\begin{enumerate}[label=\textup{(\roman*)}]
\item for every $N$, the difference in \eqref{eq:fdq-def} converges uniformly
to zero on $0\le m,n\le N$;
\item for all sufficiently small $h$, both arrays are dominated by one
summable envelope $E_{m,n}$ independent of $a,h,g,f$.
\end{enumerate}
Then \eqref{eq:fdq-def} holds.
\end{theorem}

\begin{proof}
Given $\varepsilon>0$, choose $N$ so that the sum of $E_{m,n}$ outside the
square $[0,N]^2$ is below $\varepsilon/4$.  The two tails then contribute less
than $\varepsilon/2$.  Uniform convergence on the finite square makes its
contribution less than $\varepsilon/2$ for small $h$.
\end{proof}

\begin{remark}
Pointwise convergence for each fixed $(m,n)$ is not enough.  The theorem is the
minimal dominated-convergence statement in the topology in which the response
series is actually summed.
\end{remark}

\section{Third-source totalization and the spectral three-jet}\label{sec:u3}

Let $d=(K,\varepsilon,\mathcal P)$ denote a homogeneity cutoff, a vertex or
intersection cap, and a finite atlas refinement.  For each $d$, let
\begin{equation}\label{eq:G3d}
 G_{3,d}(a)=
 \Tot_{\mathfrak S_d}\{\eta_{a,\omega}:\omega\in I_d\}
 +S_{d}^{\mathrm{seam}}(a).
\end{equation}
The totalization order $\mathfrak S_d$ records the physical assembly and is not
an arbitrary rearrangement.

Define the atom norm
\begin{equation}\label{eq:atom-norm}
 \norm{\eta_{a,\omega}}_{\mathrm{CM2}}
 =\sup_{\norm{g}\le1,\norm{f}\le1}
 \sum_{m,n\ge0}
 \abs{\pair{Q_a^n\eta_{a,\omega}Q_a^m g}{f}}.
\end{equation}

\begin{assumption}[Third-source packet]\label{ass:third}
There are numbers $b_\omega\ge0$ and moduli
$\omega_{\mathrm{vert}},\omega_{\mathrm{seam}}$ such that
\begin{align*}
 \sup_a\norm{\eta_{a,\omega}}_{\mathrm{CM2}}&\le b_\omega,
 &\sum_\omega b_\omega&<\infty,\\
 \omega_{\mathrm{vert}}(\varepsilon)&\longrightarrow0,
 &\omega_{\mathrm{seam}}(\mathcal P)&\longrightarrow0.
\end{align*}
Every finite approximant is continuous.  On a common refinement, equal UIDs
with opposite orientations cancel exactly, and no distinct occurrence is
re-keyed.
\end{assumption}

\begin{theorem}[CM2-to-$U3$ joint Cauchy theorem]\label{thm:u3}
Under \cref{ass:spectral,ass:third}, if $d_1,d_2$ refine $d$, then
\begin{equation}\label{eq:cauchy}
 \norm{G_{3,d_1}-G_{3,d_2}}_{X_0}
 \le 2\sum_{\omega\notin I_d}b_\omega
 +2\omega_{\mathrm{vert}}(\varepsilon)
 +2\omega_{\mathrm{seam}}(\mathcal P).
\end{equation}
Consequently $G_{3,d}$ converges uniformly in $a$ to a unique continuous third
source $G_3$.

If $G_j:X_r\to X_{r-j}$ and every reduced resolvent preserves its level, then
all contour-resolvent words of total parameter order at most three are bounded
and continuous.  In particular, the Riesz projection, the isolated
eigenvalue, and the reduced spectral data possess a continuous third jet.
\end{theorem}

\begin{proof}
Pass to a common refinement of $d_1$ and $d_2$.  Common atoms cancel by UID,
owner, and orientation.  The remaining difference consists of atoms outside
$I_d$, vertex-cap terms, and seam-refinement defects.  The first contribution
is bounded by twice the tail of $\sum b_\omega$ in the complete norm
\eqref{eq:atom-norm}; the other two are bounded by their defining moduli.  This
is \eqref{eq:cauchy}.  Completeness gives $G_3$, and uniform convergence gives
continuity in $a$.

For the spectral statement, differentiate the resolvent
$\cR_a(z)=(z-L_a)^{-1}$ under a fixed isolating contour.  The first three
formulas are finite sums of ordered words
\[
 \cR G_1\cR,
 \quad
 \cR G_2\cR+2\cR G_1\cR G_1\cR,
\]
\[
 \cR G_3\cR+3\cR G_2\cR G_1\cR
 +3\cR G_1\cR G_2\cR
 +6\cR G_1\cR G_1\cR G_1\cR.
\]
The graded loss exactly matches the available ladder levels.  Integrating the
words around the contour gives the Riesz projection and the corresponding
reduced data.
\end{proof}

\begin{remark}
First-order CM2 does not imply $U3$.  The product estimate must hold for the
actual derivative atoms entering the third source, and the totalization,
vertex, and seam errors must be controlled in the same complete topology.
\end{remark}

\section{An open four-branch moving-seam family}\label{sec:fourbranch}

Let $I=[-1/40,1/40]$ and define
\begin{align}\label{eq:weights}
 w_1(a)&=\frac1{10}+a,&
 w_2(a)&=\frac15+2a,&
 w_3(a)&=\frac3{10}-a,&
 w_4(a)&=\frac25-2a.
\end{align}
Let $s_0=0$, $s_i=\sum_{j\le i}w_j$, and on
$I_i(a)=[s_{i-1}(a),s_i(a))$ set
\begin{equation}\label{eq:fullbranch}
 T_a(x)=\frac{x-s_{i-1}(a)}{w_i(a)}\pmod1.
\end{equation}
This map is realized as a pinball billiard on
$\T\times[0,1]$: bottom collisions reset the vertical direction and top
collisions choose one of four winding lifts.  The three internal endpoints are
therefore physical direction seams.  Their velocities are nonzero.

The transfer operator with respect to Lebesgue probability is
\begin{equation}\label{eq:transfer-four}
 (L_a h)(y)=\sum_{i=1}^4w_i(a)
 h\bigl(s_{i-1}(a)+w_i(a)y\bigr).
\end{equation}
Lebesgue probability is invariant and $\Pi h=\int h$ is independent of $a$.

For $r\ge1$, let $X_r$ be the fixed-cut space whose restriction to $(0,1)$ is
$W^{r,1}$, with norm
\begin{equation}\label{eq:Xr}
 \norm{h}_{X_r}=\norm{h}_{L^1}+
 \sum_{j=0}^{r-1}\Var_{\T}(h^{(j)}).
\end{equation}
Circular variation includes the jump at the fixed coordinate cut.

\begin{lemma}[Uniform graded spectral contraction]\label{lem:contraction}
Let $\rho=9/20$.  For every fixed $r$ there is $C_r$ such that
\begin{equation}\label{eq:contraction}
 \norm{Q_a^n h}_{X_r}\le C_r\rho^n\norm{h}_{X_r},
 \qquad \int h=0,
\end{equation}
uniformly in $a\in I$.
\end{lemma}

\begin{proof}
Every inverse branch of $T_a^n$ is affine with image interval of length
$w_\omega\le\rho^n$, and
\[
 L_a^n h(y)=\sum_{\omega}w_\omega h(s_\omega+w_\omega y).
\]
For $0\le j\le r-1$,
\[
 \D^jL_a^nh=
 \sum_\omega w_\omega^{j+1}
 (\D^j h)(s_\omega+w_\omega y).
\]
Hence
\[
 \Var(\D^jL_a^nh)
 \le \rho^n\sum_\omega
 \Var_{I_\omega}(\D^j h)
 \le \rho^n\Var(\D^j h).
\]
The same partition estimate controls the fixed-cut jump.  For centered
functions, the circular Poincar\'e inequality controls the $L^1$ norm by the
variation.  Summing over $j$ proves \eqref{eq:contraction}.
\end{proof}

Put $v=(1,2,-1,-2)$, $A_i=\sum_{j\le i}v_j$, and
\[
 \psi_i^a(y)=s_{i-1}(a)+w_i(a)y,
 \qquad
 r_i(y)=A_{i-1}+v_i y.
\]

\begin{lemma}[Exact parameter letters]\label{lem:letters}
For every $k\ge1$,
\begin{equation}\label{eq:letter-formula}
 \partial_a^kL_a h
 =\sum_{i=1}^4\left[
 kv_i r_i^{k-1}h^{(k-1)}(\psi_i^a)
 +w_i(a)r_i^k h^{(k)}(\psi_i^a)
 \right].
\end{equation}
Moreover, for $r\ge1$,
\begin{equation}\label{eq:letter-bound}
 \norm{\partial_a^kL_a h}_{X_r}
 \le C_{r,k}\norm{h}_{X_{r+k}},
 \qquad
 \Pi\partial_a^kL_a=0.
\end{equation}
\end{lemma}

\begin{proof}
Both $w_i$ and $\psi_i^a$ are affine in $a$,
$\partial_aw_i=v_i$, and $\partial_a\psi_i^a=r_i$.  The Leibniz rule has only
the terms in which zero or one derivative hits $w_i$, giving
\eqref{eq:letter-formula}.  Composition with an affine contraction and
boundedness of the polynomials $r_i$ give \eqref{eq:letter-bound}.  The last
identity follows by differentiating mass preservation.
\end{proof}

Combining \cref{lem:contraction,lem:letters} gives, for every required level,
\begin{equation}\label{eq:model-product}
 \abs{\pair{Q_a^n(\partial_a^kL_a)Q_a^m g}{f}}
 \le C_{r,k}\rho^{m+n}
 \norm{g}_{X_{r+k}}\norm{f}_{X_r^*}.
\end{equation}
The finite ladder for third order is
\[
 X_4\xrightarrow{G_1}X_3\xrightarrow{G_1}X_2
 \xrightarrow{G_1}X_1,
\]
with the analogous $G_2$ and $G_3$ arrows.  The level $X_5$ controls the
fourth-order Taylor remainder.

\begin{lemma}[Nonzero moving source]\label{lem:nonzero}
For every $a\in I$, the first parameter letter is nonzero and has a saltus
component attached to a moving physical seam.
\end{lemma}

\begin{proof}
Choose a smooth plateau whose support meets only one moving branch endpoint
and not the fixed coordinate cut.  In fixed coordinates, differentiating the
indicator of the moving branch produces the signed boundary trace multiplied
by the nonzero seam velocity.  The two incident values of the plateau are not
equal, so the same-occurrence jump coefficient is nonzero.  No other seam can
cancel it because the UIDs differ.
\end{proof}

\begin{theorem}[Actual open CM2 and $U3$ family]\label{thm:fourbranch}
The family \eqref{eq:weights}--\eqref{eq:fullbranch} satisfies, uniformly on
$I$:
\begin{enumerate}[label=\textup{(\roman*)}]
\item a nonzero moving-seam source;
\item product-CM2 for parameter letters through order four;
\item finite difference convergence in $\ell^1(\N^2)$ through order three;
\item a finite stable occurrence atlas;
\item a continuous third spectral jet on the ladder $X_5\to\cdots\to X_1$.
\end{enumerate}
\end{theorem}

\begin{proof}
The nonzero source is \cref{lem:nonzero}; product-CM2 is
\eqref{eq:model-product}.  Taylor's formula with the fourth-order integral
remainder and the same product envelope gives the hypotheses of
\cref{thm:fdq}.  The atlas consists of four branches and three physical seams,
so no infinite refinement is required.  Finally, \cref{thm:u3} applies and the
displayed ladder makes every third-order resolvent word well typed.
\end{proof}

\section{A nonconjugate specular Sinai response channel}\label{sec:specular}

Let $Q_a$ be a $C^{J+6}$ family of periodic planar finite-horizon dispersing
billiards.  All scatterers are fixed except one circle of radius $R_1+a$.
Assume that at $a=0$ the normal two-cycle between this circle and a fixed
circle is the unique shortest regular period-two orbit, separated from all
other period-two lengths by a positive gap.  These conditions persist for
small $a$.

Let $d(a)=d_0-a$, $\kappa_1(a)=(R_1+a)^{-1}$ and
$\kappa_2=R_2^{-1}$.  In paraxial coordinates the return matrix is
\[
 M(a)=S(\kappa_1(a))F(d(a))S(\kappa_2)F(d(a)),
\]
where
\[
 F(d)=\begin{pmatrix}1&d\\0&1\end{pmatrix},
 \qquad
 S(\kappa)=\begin{pmatrix}1&0\\2\kappa&1\end{pmatrix}.
\]
A direct multiplication gives
\begin{equation}\label{eq:trace}
 \operatorname{tr}M(a)=2+4d(a)(\kappa_1(a)+\kappa_2)
 +4d(a)^2\kappa_1(a)\kappa_2.
\end{equation}
Its derivative at zero is strictly negative.  The continuous-time period is
$2d(a)$, which also varies.  The isolation gap prevents a conjugacy from
permuting this orbit with another one.  Thus the flow is not time-preservingly
conjugate, and the collision maps are not locally $C^1$ conjugate.

Use normalized arclength coordinates on every collision component.  If
$\ell_i(a)$ is the physical perimeter and $\ell(a)=\sum_i\ell_i(a)$, then with
respect to $dm=dx\,d\varphi$ the invariant density is
\begin{equation}\label{eq:sinai-density}
 \rho_a(i,x,\varphi)=
 \frac{\ell_i(a)\cos\varphi}{2\ell(a)}.
\end{equation}
It is real analytic in $a$.

\begin{theorem}[Assembled specular radial response]\label{thm:sinai-u3}
For every finite $J$ allowed by the boundary smoothness and the finite-order
source chart, differentiating $L_a\rho_a=\rho_a$ gives complete sources
\begin{equation}\label{eq:assembled-source}
 S_j=\sum_{k=1}^j{j\choose k}L_0^{[k]}\rho_0^{(j-k)}
 =(I-L_0)\rho_0^{(j)},
 \qquad 1\le j\le J.
\end{equation}
Every $S_j$ is a centered strong density after complete physical assembly and
\begin{equation}\label{eq:assembled-recovery}
 R_0S_j=\rho_0^{(j)}.
\end{equation}
In particular the nonconjugate radial family has an actual invariant-projector
$U3$ channel.

If
\[
 \psi_a=c(a)+g_a-g_a\circ T_a,
\]
then the twisted transfer operator satisfies
\begin{equation}\label{eq:gauge}
 L_{a,q}=e^{qc(a)}M_{e^{-qg_a}}L_aM_{e^{qg_a}}.
\end{equation}
Consequently the leading eigendata in every such cohomological channel have all
mixed derivatives allowed by $a\mapsto(c(a),g_a,\rho_a)$.
\end{theorem}

\begin{proof}
The explicit formula \eqref{eq:sinai-density} gives the density jets.
Differentiate invariance in the distributional source module.  The Leibniz
rule gives the left side of \eqref{eq:assembled-source}; rearrangement gives
the smooth right side.  Normalization implies
$\int\rho_0^{(j)}dm=0$, so the reduced resolvent gives
\eqref{eq:assembled-recovery}.  For the twist, use
$L_a((u\circ T_a)h)=uL_ah$ to obtain \eqref{eq:gauge}; the eigendata follow by
explicit conjugation.
\end{proof}

\begin{remark}[Scope]
The theorem concerns the invariant projector and exact cohomological twists.
It does not assert differentiability of the full reduced resolvent for an
arbitrary noncoboundary observable.  Such a statement requires the bilateral
third-source packet of the preceding sections.
\end{remark}

\section{Maximality and obstruction theorems}\label{sec:maximality}

The positive theorem cannot be reduced to a finite list of local geometry
checks.

\begin{proposition}[Finite prefixes do not determine CM2]\label{prop:prefix}
Fix any finite occurrence registry and any finite rectangle of the arrays
\eqref{eq:atom}.  There are two completions agreeing on this data, one with a
summable product tail and one for which the absolute double sum diverges.
\end{proposition}

\begin{proof}
For the first completion, set every unrecorded atom to zero.  For the second,
add disjoint formal occurrences outside the finite prefix with positive matrix
elements $(m+n+1)^{-1}$ on a sparse diagonal family chosen so that each finite
rectangle is unchanged but the total sum diverges.  Both completions satisfy
the same finite records.  Thus no finite-prefix certificate can imply a global
$\ell^1(\N^2)$ conclusion.
\end{proof}

A geometric version uses the trace defect map.  Let $\mathcal V$ be a Banach
space of localized primitive jets near a regular moving face, and let
\[
 \mathfrak D:\mathcal V\to C^r(J)
\]
assign the unrecovered terminal trace on a compact face subarc $J$.

\begin{theorem}[Maximality of the recovery condition]\label{thm:maximality}
Suppose $\mathfrak D$ has a continuous right inverse.  Then
$\ker\mathfrak D$ is complemented and has infinite codimension.  Its complement
is open and dense after restricting to any affine source class on which
$\mathfrak D$ is nonconstant.  Hence local finite-horizon geometry, compactness,
and finitely many Ward identities do not imply recovery of arbitrary moving
face sources.  A valid general theorem must contain either a source-specific
complete assembly or an explicit bilateral product-tail packet.
\end{theorem}

\begin{proof}
If $S:C^r(J)\to\mathcal V$ is a right inverse, then
$P=\Id-S\mathfrak D$ is a bounded projection onto $\ker\mathfrak D$.  Since
$C^r(J)$ is infinite dimensional, the quotient
$\mathcal V/\ker\mathfrak D$ has infinite dimension.  On an affine source
class, the passing set is an affine translate of the kernel.  A proper closed
subspace has empty interior, so its complement is open and dense.  The same
argument applies at successive jet orders to the corresponding affine defect
maps.
\end{proof}

\section{Consequences and outlook}

The theory isolates three independent ingredients of higher response with
moving singularities:
\begin{enumerate}[label=\textup{(\arabic*)}]
\item bilateral decay in the two time variables;
\item absolute totalization of the physical source atoms;
\item a graded scale on which every Kato word is composable.
\end{enumerate}
The crossed-rate theorem handles the first, the joint Cauchy theorem the
second, and the loss ladder the third.  The open four-branch family proves that
all three can hold with a nonzero moving source.  The radial Sinai family shows
that complete physical assemblies can also close in a genuinely nonconjugate
specular system.  The maximality theorem explains why the corresponding
hypotheses are structural rather than technical artifacts.

The companion paper develops pressure, suspension resolvents, and physical
diffusion from the third spectral jet established here.

\section*{Acknowledgements and disclosure}
The author is responsible for every mathematical statement and citation in the
circulated version.  This referee draft was assembled with computational and
language-model assistance; all proofs require human verification before
submission.

\bibliographystyle{amsplain}
\bibliography{references}

\end{document}
